% 本文件由 LaTeX 排版 Agent 根据有效 translation.json 自动生成。
% author=Rufinus; work_id=2712; title=诸序言

\chapter{诸序言}

% structure: p0001 source=h1 target=omitted reason=page-title-already-rendered-by-parent

% structure: p0002 source=h2 target=section reason=relative-heading-rank; base=h2

\section{圣克肋孟\WorkTitle{认语录}诸书序言}

\Emph{致高登提主教}\par

（关于本作品的写作场合与日期，见绪论，第412页。）\par

亲爱的\Person{高登提}，你品格如此刚强，堪称我们教师的杰出荣耀；或者我宁愿说，你蒙受了圣神的恩宠如此丰沛，以致连你日常交谈或你在教会中讲道时所说的话，都应笔录下来，传予后世作为教训。但我远没有那么敏捷，天资浅薄，年迈已使我迟钝缓慢；而这工作无非是偿还债务，因我所敬重的贞女\Person{西尔维亚}曾命令我，正如你现在以继承权所做的那样，将克肋孟的著作译成我们的语言。正是她曾要求我，如同你现在凭继承权所要求的，将克肋孟的著作译成我们的语言。这笔债终被偿还，虽历经多次拖延。这是我从未从希腊人的图书馆中携获的战利品之一部分，在我看来并非微不足道，我正为同胞的使用和利益而收集它们。我没有自己的食物带给他们，必须从国外输入滋养品。不过，外来的货物往往显得更甜美，有时也确实更有用。此外，凡能治愈我们身体、防御疾病或解毒之物，几乎都来自国外。\Place{犹太}给我们送来香树脂的蒸馏液，\Place{克里特}送来白藓叶，\Place{阿拉伯}送来芳香之花，\Place{印度}送来甘松之实。这些货物来到我们这里，无疑不如我们田地所产的完美，但它们完好地保存了其芳香与疗愈之力。因此，我的如同我自己灵魂的朋友，我将重返罗马的克肋孟呈献给你。我将他打扮成拉丁装束呈献给你。倘若他的辞采所呈现的面貌不如可能那般明亮，请勿诧异。只要意义相同，便无差别。\par

那么，这些便是外来的货物，我以极大的劳力输入；我仍需看我的同胞们是否会感激那个为他们带来\OriginalTerm{战利品}{spolia}的人，他正以我们语言的钥匙开启一个迄今隐藏的宝库，虽然他这样做是出于最大的善意。我只信赖天主会眷顾你的善意，使我的礼物在任何地方都不遭恶眼与嫉妒之视；使我们不致目睹那极其怪异的现象：受赠者表达恶意，而被取走者却毫不吝惜地交出。你既已读过此作品的希腊文，理当向他人指出我翻译的意图——除非你确实觉得在某些方面我没有遵循正确的翻译方法。我相信你很清楚，克肋孟这部\WorkTitle{认语录}有两个希腊文版本；有两套书，它们在少数情况互有差异，但叙述主体相同。例如，作品最后一部分，即记述\Person{西满魔师}变形的那部分，存在于其中一个版本，而在另一个版本中完全缺失。另一方面，有一些内容，例如关于无始与受生的天主的论述，以及少数其他内容，虽然在两个版本中都有，但至少可以说，超出我的理解；我宁愿留待他人处理，也不愿以不足的方式呈现它们。至于其余部分，我苦心不偏离其意义或措辞，即使最微小的程度也不；这虽使表述不够华丽，却使之更为忠实。\par

有一封信，同一克肋孟致\Person{雅各伯}——主的兄弟，在其中叙述了\Person{伯多禄}的死，并说他已立\Person{雅各伯}为他的继承人，作为教会的治理者和教师；此外还包含了一整套教会治理方案。我没有将这封信置于作品之前，既因它在时间上较晚，也因我之前已将其翻译并出版。然而，有一点若我将此信的翻译置于本作品内，或许会显得与事实不符，但我不认为它涉及不可能性。那就是：\Person{利诺}和\Person{克莱图}在克肋孟之前已是罗马城的主教。那么，有些人问，克肋孟在致雅各伯的信中怎能说伯多禄将教会教师的职位传给了他呢？按我的理解，这点的解释如下：利诺和克莱图无疑在克肋孟之前是罗马城的主教，但这发生在伯多禄在世之时；也就是说，他们负责主教工作，而伯多禄履行宗徒的职责。已知他在\Place{凯撒勒雅}也做了同样的事；因为在那里，虽然他自己在场，却任命了\Person{匝凯}为主教，留在他身边。如此我们可看出两方面如何可能都是真实的：即他们如何列在克肋孟之前的主教名单上，而克肋孟在伯多禄死后又如何成为他在教师席位上的继承人。但现在我们该注意克肋孟本人叙述的开端，这是他致主之兄弟雅各伯的。\par

% structure: p0008 source=h2 target=section reason=relative-heading-rank; base=h2

\section{第二卷序言}

\Emph{\Person{鲁菲努}，从他对前书的序言可以看出，认为只阐释对\Person{犹大}的祝福而不及其兄弟是不令人满意的。因此，\Person{保利努}借他们共同的朋友\Person{塞雷亚利斯}前往罗马之机，致函\Person{鲁菲努}，劝其阐释其余的祝福。}\par

保利努致敬其兄弟鲁菲努，并致一切美好祝愿。\par

1. 虽然我们的儿子\Person{塞雷亚利}（Cerealis）对我说，他如今返回\Person{圣伯多禄}那里时，是否能探望您并不确定；然而在我看来，如果我不藉着他——这位您和我共同拥有的人——写信给您，那将是我应受责备的事，也会使您烦恼。我认为，宁可因他不来探望您而损失一些信纸，也不可因他空手前来而失去您的信任，我相信后一种情况会发生。因此我将这封信——我不说是托付给机遇，而是托付给信德——因为我相信，\Emph{主}必会引导我们的儿子和我的信的道路到您那里；因为凡渴望善的人，一切都会转为善；他确实渴望您，正如一个懂得从与您的交往中获得益处的人所应渴望的那样。我相信，他对这一善事的渴望不会落空，这是根据他的信德和孝爱；所以我有信心，他必能到达您那里并与您同住，我也将看到\Emph{主}的救恩加倍地赐予您，因为您将在他身上获得一位好儿子、学生和助手的加入，而他将从您那里找到一位父亲和老师，教导一切从\Emph{主}那里赐予他的美好事物，\Emph{主}将加强他祈祷的效能和力量，赐予属神恩宠的能力。至于我自己，虽然我有把握，当您返回东方时，您不会不愿意前来探望我；但我的罪使我害怕，恐怕巴比伦的女儿会使您远离我。因此，我以热切的渴望向\Emph{主}祈求，求祂不按我的功过，而按我的渴望赐予我，并引领您的行程在平安的道路上到我这里来；因为不在这条路上行走的人，是被弃绝和定罪的，无法真正渴望您的临在。\par

2. 但现在要谈我信中的正事了。我以惯常的恳切态度要求您——即使深更半夜我也习惯这样敲您的门，因为担心遭到拒绝而不得不保持谦卑祈求者的姿态——再一次向我施恩，阐释那对十二位圣祖的祝福。您已经开始了与\Person{犹大}有关的预言，并按照诫命作了三层面的解释。我现在请求您阐释有关\Person{犹大}每个儿子的预言：这样，我本人就能藉着您而拥有真理，并藉着您的帮助获得将归于我的恩宠和赞美；因为我将能回答那些认为应当就这段圣经的疑难向我请教的人，不是凭我自身理解所出的愚昧话语，而是凭您所启示的神圣真理。\par

————————————\par

\Person{鲁菲努斯}（Rufinus）虽然此时正忙于自己的大作——翻译\Person{庞菲鲁斯}（Pamphilus）为\Person{奥利振}（Origen）的辩护，以及\Person{奥利振}的\OriginalTerm{论原理}{\GreekText{Περὶ} \GreekText{᾽Αρχῶν}}——并且即将动身前往罗马，但他毫不迟疑地着手完成了\Person{保利努斯}（Paulinus）所要求的著作，并随附以下书信寄给他。\par

\Person{鲁菲努斯}致其兄弟\Person{保利努斯}——天主之人——并致一切良好祝愿。\par

1. 虽然我们共同的儿子\Person{塞雷亚利}没有来探望我，但他知道若延迟我收到您的信会给我带来多大痛苦，于是将信转寄给了我。在阅读时，我像往常一样，对您的思念不断加深；但在近结尾处我发现了一个请求——我曾多次请您免去这个请求——即您要求我就圣经段落的解释问题写一些答复。我以为，通过我一再寄给您的那些著作（它们证明了我的无知和言语粗陋），我本应使您放弃这些问题的。\par

2. 但既然您仍不厌倦命令我，我立即尽我所能，在我按您要求所写的关于\Person{犹大}祝福的注释之外，又加上了其余十一位圣祖的注释。我像那比喻中两个儿子的第一个人。我想这样最能承行父亲的旨意：虽然当他吩咐我到葡萄园去时，我说\Quote{我不去}，但过后我还是去了。我承认，以如此微薄的能力承担如此重任，不免有些鲁莽，但我恭敬地对您说，这必须最公正地归咎于您，因为您对我过分的爱，使您看不到我的知识——如同其他德行——只是渺小的。我在四旬期内，寓居在\Place{松林隐修院}时写成了这部著作，并且是为您而写。但我觉得无法向那里的弟兄们隐藏这部拙作；他们认为，一件蒙您赞同的事必定非常重要，便强求我允许他们为自己抄写。这样，当您向我索取食物为自己时，您也滋润了他人。再见，平安，我最亲爱的兄弟，最真诚的天主敬拜者，一个没有诡诈的以色列人。我恳求您——充满天主恩宠的人——仍然记念我。\par

% structure: p0018 source=h2 target=section reason=relative-heading-rank; base=h2

\section{致其翻译的\Person{庞菲鲁斯}\WorkTitle{为奥利振辩护}序言}

写于\Place{松林隐修院}，公元三九七年。\par

当鲁菲努斯居住在皮内屯时，一位名叫玛加略的基督徒向他寻求建议和帮助。此人正与\Quote{数术家}\OriginalTerm{数术家}{Mathematici}一派人进行论战，这些人已背离了其得名的科学研究，转向占星术与宿命论信仰。玛加略听闻奥利振在基督宗教思辨领域的卓绝成就，热切渴望了解其著作，却因不通希腊文而无法实现。他向鲁菲努斯宣称，自己曾梦见一艘满载东方货物的船只抵达意大利，且有人向他启示，这艘船将载有获得他所渴望之知识的方法。鲁菲努斯的到来在他看来正是梦境的应验，于是他恳切请求鲁菲努斯将希腊学问的宝藏分赐于他，尤其希望他将奥利振那部伟大的思辨著作\WorkTitle{\OriginalTerm{论首要原理}{\GreekText{Περὶ} \GreekText{᾽Αρχῶν}}}即\WorkTitle{论首要原理}翻译出来。鲁菲努斯起初犹豫，因为他深知世人对奥利振怀有强烈偏见，尤其在西方，人们视其为异端，尽管其著作在当地鲜为人知。然而，他终究屈服于玛加略的恳求；但为防范被指为异端，他着手撰写了三部预备性作品：首先，他翻译了殉道者庞斐禄为奥利振所写的\WorkTitle{辩护辞}；其次，他撰写了关于异端者窜改奥利振著作的短论；第三，在翻译\WorkTitle{\OriginalTerm{论首要原理}{\GreekText{Περὶ} \GreekText{᾽Αρχῶν}}}时，他为之撰写了一篇详尽的序言，以证明自己翻译此作的合理性。所有这些文献都引发了激烈的争论，体现在耶柔米致罗马友人的信函以及本卷所译的鲁菲努斯与耶柔米的\WorkTitle{辩护辞}中。\par

庞斐禄为奥利振所作的\WorkTitle{辩护辞}是他与教会史家凯撒勒雅的欧瑟伯合作撰写的一部著作的第六卷。庞斐禄是一位伟大的藏书家与博学之士，但欧瑟伯是主要的执笔者。庞斐禄在最后一次迫害（即加莱里乌斯治下的迫害）中殉道；后来欧瑟伯因涉嫌\OriginalTerm{亚略主义}{Arianism}而遭猜疑，于是那些不喜欢奥利振的人试图将庞斐禄与这部著作完全脱离干系。然而，尽管耶柔米激烈抗议，似乎仍无理由怀疑庞斐禄全程参与了该著作，并且确实撰写了第六卷。这部\WorkTitle{辩护辞}的翻译是最先完成的，并附有一篇序言，其内容如下：\par

你渴望认识真理，极可爱的\Person{玛加略}，所以请我为你作一件事，这件事将使你获得认识真理的福份；但这件事却会替我招来那些人的极大愤怒，他们总以为，只要有人不认为\Person{奥力根}是坏蛋，就是得罪了他们。的确，你问的不是我本人对他的看法，而是\Person{圣殉道者庞斐禄}的看法；你并要求将据说他用希腊文所写的辩护书翻译给你，译成拉丁文。然而，我不怀疑，即使我用他人的话来为他辩护，也会有人认为我冒犯了他们。我恳请他们不要存有傲慢和偏见的心；既然我们都要站在基督的审判台前，就请他们不要拒绝听真理的话，免得因无知而做错事。让他们想想，用诬告来伤害软弱弟兄们的良心，就是得罪基督。因此，不要听信控告者，也不要从第三者那里打听别人的信仰，尤其是当有机会亲自、直接地认识，并且各人的信仰实质与品质是由他自己宣认而得知的时候。因为圣经说：\BibleQuote{人心里相信，就得成义；口里宣认，就得救恩}。又说：\BibleQuote{凭你的话，你要被成义；凭你的话，你也要被定罪}。\Person{奥力根}在各处圣经中的见解，在本作品中清楚陈述。至于为何我们会发现某些地方他自相矛盾，将在所附的简短文件中加以解释。至于我自己，我坚持从圣教父们传下来的教导，即天主圣三同为永远，同属一个本性、德能与实体；天主子在这末世降生成人，为我们的过犯而受苦，并在他受苦所在的肉身中从死者中复活，借此将复活的希望赐给全人类。当我们谈论肉身的复活时，我们并非如某些人恶意诽谤的那样使用遁词；我们相信就是要复活的这同一个肉身，即我们现在所生活的肉身，不是另一种肉身，也不是与此肉身不同的另一个身体。当我们说身体复活时，我们用的是宗徒的话，因为他自己用了这个词；当我们说肉身时，我们的宣认是按信经。认为人的身体不同于肉身，这是恶意的荒谬杜撰。然而，无论我们按着共同的信仰称那将要复活者为肉身，或者按着宗徒称它为身体，我们都必须相信，按着宗徒清楚的话，那将要复活的，是在大能和光荣中复活；它将复活为不朽的、属神的身体；因为朽坏不能承受不朽。我们必须坚持那将要有的身体或肉身的卓越性；但在此条件下，我们必须坚持肉身的复活是完全而完整的；一方面要维护肉身的同一性，另一方面又不能减损不朽和属神身体的光荣与尊严。因为圣经是这样说的。这就是在耶路撒冷的可敬主教\Person{若望}所宣讲的；这就是我们与他一同宣认和坚持的。若有人相信或教导与此不同，或暗示我们的信仰与此不同，就让他受诅咒。愿任何愿意知道的人，将这当作我们信仰的记录。我们所读和所做的，都符合这信仰的陈述；我们跟随宗徒的话：\BibleQuote{考验一切，持守善的，远避各种邪恶}。凡按此准则而行的人，愿平安归于他们，也归于天主的以色列。\par

% structure: p0023 source=h2 target=section reason=relative-heading-rank; base=h2

\section{\Person{鲁菲诺}为\Person{殉道者庞斐禄}的\WorkTitle{为\Person{奥力根}辩护}所写的跋}

\Emph{又题为\WorkTitle{论奥力根著作的窜改}}\par

致住在\Place{毕内屯}的\Person{玛加略}，公元397年。\par

下一篇作品是与\Person{庞斐禄}的\WorkTitle{辩护}同时发出的。\Person{鲁菲诺}相信，\Person{奥力根}的著作已被异端者窜改，以致把他的论断变成支持他们自己观点的论据。因此，在他翻译\WorkTitle{\OriginalTerm{论首要原理}{\GreekText{Περὶ} \GreekText{᾽Αρχῶν}}}时，他更改了许多在\Person{奥力根}常见手稿中带有异端含义的段落，以使该著作与自身一致，并与\Person{奥力根}其他著作中表达的正统观点保持一致。这种做法是否正当或诚实，必须从随后发生的争论来判断；但应记住：首先，文学精确性和尽责的标准在那时不同于我们这时代；其次，当一切都依赖抄写员时，抄本中就会有无穷的变化，无论是由于疏忽、无知还是欺诈；第三，\Person{鲁菲诺}所采纳的原则，正是他的伟大对手\Person{热罗尼莫}在他的\WorkTitle{论翻译的最佳方式}以及他的\WorkTitle{致维吉兰修书}（第66和第61封）中所承认的原则。\par

我将圣殉道者\Person{庞斐禄}的著作\WorkTitle{为奥力振辩护}从希腊文译为拉丁文，这是我在前一卷中按照我的能力和事情的需要所做的。其目的是：我希望你通过充分的信息知道，在他的著作中已阐述的信德准则正是我们必须接受和坚持的；因为清楚地显示公教观点包含在其中。然而，你必须承认，他的书中有一些东西不仅与此不同，而且在某些情况下甚至与之相悖；这些东西是我们真理准则所不认可的，我们既不能接受也不能赞同。至于这原因，有一个普遍持有的观点传到了我这里，我希望你和那些渴望知道真理的人都能完全知道，因为也有可能一些先前受挑剔之爱驱使的人，当事实摆在他们面前时，会默认这件事的真理和理性；因为有些人似乎决心相信世界上任何东西都是真实的，而不相信那些使他们失去挑剔机会的东西。我认为，人们必然感到完全不可能：一个如此博学、如此智慧的人，一个连他的控告者也完全可以承认既非愚蠢也非疯狂的人，竟然会写出与他自己的观点相悖和冲突的东西。但即使假设这在某种程度上可能发生；假设，如有些人也许说过的，他在晚年可能忘记了他早年所写的东西，并发表了与先前观点不一致的断言；我们又当如何处理这样一个事实：我们有时在完全相同的段落中，甚至如我所说，几乎在连续的句子里，发现有表达相反意见的从句插入？我们怎能相信，在同一部著作、同一卷书中，甚至如我所说，有时在紧接的段落中，一个人会忘记自己的观点？例如，当他刚刚说过，在全部圣经中找不到任何一处说圣神是受造或被造，他可能立刻又补充说圣神与其余受造物一同受造？或者，同样一个人，明明说圣父与圣子是同一实体，或在希腊文中称为\OriginalTerm{同质}{\GreekText{ὁμοούσιον}}，却在下一句说他属于另一实体，是受造物，而就在此前不久他还描述他为出自天主圣父的本性而生？或者，关于肉身的复活，他如此清楚地宣告，正是肉身的本性连同天主的圣言升了天，并在那里向天上的权能者显现，向他们呈现一个可供敬拜的新形象，那么他，我问你，有可能转过头来说这肉身不会得救吗？即使是一个失去理智、头脑不清的人，也不可能发生这样的事。因此，我将尽可能简短地指出这是如何发生的。异端者能行任何暴行，他们毫无悔意，也毫无顾忌：我们不得不从他们屡次被证实的大胆妄为中认识到这一点。并且，正如他们的父亲魔鬼从一开始就以伪造天主的话、扭曲其真正含义为目的，并狡猾地将他自己有毒的思想插入其中，同样地，他也将同样的技艺作为遗产留给了这些后继者。因此，当天主对\Person{亚当}说：\Quote{园中所有树上的果子你都可以吃；}他（魔鬼）想要欺骗\Person{厄娃}，便插进了一个音节，藉此将天主慷慨允许吃所有果子的范围缩到最小。他说：\Quote{天主真的说过，\InnerQuote{你们不可吃园中\Emph{任何}树上的果子}吗？}这样，通过暗示天主的命令是严厉的，他更容易说服她违反诫命。异端者效法了他们父亲的榜样，他们老师的诡计。每当他们在古代著名作家的作品中发现有关天主荣耀的如此充实而忠实的论述，以至于每个信徒都能从中获益和受教时，他们便毫无顾忌地将自己错误教义的有毒污点注入那些著作中；他们这样做，要么是插入作者没有说过的话，要么是通过增补篡改作者说过的话，以便他们有毒的异端能更容易地借着所有最博学、最著名的教会作家的名义被主张和认可；他们意在让人觉得那些知名且正统的人也持有与他们相同的观点。我们在希腊作家身上有最清晰的证据，这种对书籍的篡改在许多古人身上都可以找到；但引用少数见证就足够了，以便更容易理解\Person{奥力振}的著作遭遇了什么。\par

\Person{克莱孟}，宗徒的门徒，继宗徒之后成为罗马教会的主教，是位殉道者。他写了一部作品，希腊文称为\OriginalTerm{认主}{\GreekText{Αναγνωρισμός}}，拉丁文称为\WorkTitle{认主}。在这部书中，他再三以宗徒\Person{伯多禄}的名义阐发一种看似真正宗徒性的教义；然而在某些段落中，\Person{欧诺缪斯}的异端被如此引入，以至你会以为自己在聆听\Person{欧诺缪斯}本人的论证，断言天主子是从无中创造的。接着又引入了另一种窜改方法，企图表明魔鬼和其他邪魔的本性并非来自他们意志和意图的邪恶，而是出于他们受造时一种异常且特殊的品质，尽管他在所有其他地方都教导说，每个有理智的受造物都赋有自由意志的能力。他的书中还插入了一些教会信经所不容的内容。那么我问，我们对这些事情应当怎样想？我们是否应当相信，一位宗徒性的人物，不，几乎就是一位宗徒（因为他所写的是宗徒们所讲的），一位宗徒\Person{保禄}曾以这些话为他作证：\BibleQuote{与\Person{克莱孟}和其他同工，他们的名字都在生命册上}的人，竟然写出了与生命册相矛盾的话？或者，我们是否应当像前面说过的那样说，是那些乖僻的人，为了借圣人的名字为自己的异端获得权威以便更容易被接受，而窜改了这些内容，使我们无法相信真正的作者曾想过或写过这些东西？\par

还有另一位\Person{克莱孟}，\Place{亚历山大}的司铎，该教会的教师，在他的几乎所有著作中都描述三位具有同一的荣耀和永恒：然而我们有时在他的书中发现一些段落，其中他说圣子是天主的受造物。像他这样伟大、在所有方面都如此正统、如此博学的人，难道会持有相互矛盾的观点，或者留下关于天主的见解，这些见解不仅相信，就连听一听都是不敬吗？\par

还有，\Place{亚历山大}的主教\Person{狄约尼削}，是教会信仰最博学的维护者，在无数段落中捍卫天主圣三的合一与永恒，如此热切以致一些见识较浅的人以为他持有\Person{撒柏流}的观点；然而在他反对\Person{撒柏流}异端的著作中，却插入了一些内容，以致亚略人试图借他的权威为自己辩护，为此圣主教\Person{亚大纳削}不得不为他的作品写一篇辩护词，因为他确信\Person{狄约尼削}不可能持有奇怪的见解或写下自相矛盾的话，而认为这些东西是被心怀恶意的人曲解了。\par

这一意见是由事实本身的力量引导我们形成的，在这些极为可敬的教会教师身上，我们发现，我断言，无法相信这些曾一再支持教会信仰的可敬人士会在某些特定点上持有与自己矛盾的观点。然而，对于\Person{奥力振}，正如我上面所说，在他身上，如同在其他那些人身上一样，可以找到某些表述上的差异；但仅仅以我们对那些享有正统声誉之人的想法或感受来对待他，是不够的；若不是他自己的言语和著作证明了这一事实——他曾在其中恳切地抱怨这一情况——类似指控也无法用类似的借口来应对。当他尚在世肉身中、仍有感觉和视觉时，他的书籍和论著遭到窜改或出现伪造版本，他所遭受的苦难，我们可以从他写给\Place{亚历山大}某些亲密朋友的信中清楚地了解；由此你将明白，为什么他的著作中会出现一些自相矛盾之处。\par

有些喜欢指责邻人的人，将亵渎的罪名加在我们和我们的教导上，尽管他们从未从我们口中听到过任何此类言论。他们应当留意自己，如何拒绝注意那庄严的训诫——说\BibleQuote{\BibleRef{格前6:10} 辱骂者不得继承天主的国}——当他们宣称我主张那邪恶与丧亡之父、并从天主的国中被驱逐者的父，即魔鬼，将要得救时，这是任何人即使丧失理智、明显疯狂也不能说的话。然而，我以为，我的教导被对手歪曲，并像宗徒保禄的书信一样被败坏和篡改，并不足为奇。我们知道，有些人冒保禄之名编造了一封伪书信，以便扰乱得撒洛尼人，仿佛主的日子临近了，从而欺骗他们。正是为了那封伪书信，他在\WorkTitle{得撒洛尼后书}中写了这些话：\BibleQuote{\BibleRef{得后2:1-3} 弟兄们，我们因我们的主耶稣基督的来临，和我们聚集到他面前，请求你们，不要因着某种神恩、言语或像出于我们的书信，迅速动摇心思，也不可惊慌，好像主的日子迫近了。无论如何，不要让人欺骗你们。}我看出同样的事也正发生在我们身上。某个异端推动者，在我们于许多人面前进行讨论、并做了笔录之后，从笔录者那里取得了文件，随意增删，按自己认为正确的更改，并将之公开，仿佛是我的作品，却对他自己插入的措辞得意地指摘。巴勒斯坦的弟兄们对此愤慨，派人到雅典来找我，以获取我作品的真本。那时我甚至从未重读或修订过那作品；它已被完全忽略和丢弃，几乎找不到。然而，我还是寄去了；——天主作证，我说的是实话——当我遇到那个篡改作品的人本人，并责备他如此行时，他回答，仿佛是在给我满意：\Quote{我这样做，是因为我想改进那篇论文，清除它的错误。}他对我论文的\Quote{清除}是怎样的呢？正如马西昂或他之后的阿培肋对\WorkTitle{福音}和\WorkTitle{宗徒著作}所行的清除。他们颠覆了\WorkTitle{圣经}的真本；而此人同样先删除了我做的真实陈述，然后插入虚假的东西，以提供控告我的理由。但是，虽然那些胆敢这样做的人是不虔敬和异端之人，但那些相信这种对我们的控告的人，将无法逃避天主的审判。另外还有不少人，出于给教会制造混乱的愿望而这样做。最近，某个异端者在厄弗所见过我，却拒绝见我，在我面前未开口，因某种原因回避了；之后他按自己的幻想写了一篇论文，部分是我的，部分是他自己的，并寄给他在各地的门徒：我知道它传到了罗马的人那里，也不怀疑它到了其他人那里。他在安提约基雅也以同样的鲁莽行事，在我去那里之前：他带去的论文落到了我们许多朋友手中。但我到达后，在许多人面前责备了他；当他坚持，毫无羞耻地无耻辩护其伪造时，我要求将书拿到我们中间，好让弟兄们辨认我的说话方式——他们自然知道我惯常强调的要点和所用的教导方法。然而，他不敢把书拿来，他的说法被众人驳斥，他自己被证明是伪造者，于是弟兄们得了一个教训，不要听信这种控告。因此，如果有人愿意信任我——我在天主面前说话——就请相信我对那些被虚假插入我书信中的事情所说的话。但若有人拒绝相信我，并选择诽谤我，他所伤害的不是我：他将在天主面前被指控为假证人，因为他要么作假证反对邻人，要么相信那些作假证的人。\par

这就是他尚在世时，当他还能察觉自己著作中被篡改和伪造之处时所发出的抱怨。他还有另一封信，我记得其中读到过对伪造其作品的抱怨；但我手边没有副本，否则我可以再引用一个直接来自他本人的、关于他真诚和真实的第二份证言，以补充我已引用的那些。但我认为，对于那并非出于争辩和诽谤，而是出于热爱真理而听此言的人，我已说得足够多。我已证明并证实，在我所提及、其正统性毋庸置疑的圣人们的事例中，当一本书的内容整体看似正确时，其中任何与教会信仰相悖之处，更应相信是被异端者插入，而非作者所写：我认为要求对奥力振也应相信同样的事并非荒谬，不仅因为论证相似，而且因为他本人在我从其著作中引出的抱怨中作了见证：否则我们必须相信，像一个愚蠢或疯狂的人，他写了与自己矛盾的话。\par

关于异端者可能以所推测的暴烈方式行事的可能性，我们很容易相信他们有这样的邪恶。他们已提供了一个样本，使目前的情况可信，即他们甚至无法阻止自己亵渎的手触及福音的神圣经文。任何人若想看他们在宗徒大事录或他们的书信中如何行事，如何玷污和啃噬它们，如何以各种方式亵渎它们，有时添加表达其邪恶教义的词语，有时删除反对的真理，若他阅读\Person{戴尔都良}反驳\Person{马尔西翁}的著作，就会最充分地理解。他们竟敢篡改天主我们救主的言语，那么篡改\Person{奥利振}的著作也就不足为奇。诚然，有些人可能会因异端的不同而不同意我所说的；因为是一种异端的支持者篡改了福音，而另一种异端则是这些被我们断言插入\Person{奥利振}著作中的段落所针对的。让那些有如此疑虑的人想一想：正如在众圣人内居住着天主的同一圣神（因为宗徒说：\BibleQuote{先知的神魂隶属于先知}；又说：\BibleQuote{我们都曾喝过那同一的圣神}），同样，在众异端者内也居住着魔鬼的同一邪灵，他时刻教导他们同样或类似的邪恶。\par

然而，有些人可能认为我们给出的例子说服力较弱，因为涉及的是希腊作者；因此，尽管我正在为之辩护的是位希腊作者，但由于可以说托付给这个论辩的是拉丁语，而且你们求我热心在拉丁人面前为这些人辩护，并展示他们的作品因诽谤性的翻译而遭受何等创伤，那么证明同类事情也发生在拉丁和希腊作者身上，以及那些因圣德品格而被认可的人，因自己作品的篡改而遭受诽谤风暴，将是令人满意的。我将讲述仍记忆犹新的事情，以便我论点的明显可信性毫无欠缺，其真理向所有人敞开。\par

\Place{普瓦捷}的主教\Person{依拉利}是公教教义的信仰者，并写了一部非常完整的训导著作，旨在使那些签署了背信弃义的\Place{阿里米诺}信经的人从错误中回头。这本书落入了他的对手和恶意者之手，正如一些人所说，无论是通过贿赂他的秘书，还是其他原因。他对此一无所知：但书被他们篡改了，而这位圣德之人始终完全不知情，以致当他的敌人在主教会议上开始控告他持异端，认为他持有他们知道已腐败地插入其手稿中的内容时，他本人要求拿出他的书作为他信仰的证据。书从家中取来，发现充满了他否认的内容：但这导致他被绝罚，并被排除在教务会议之外。然而，在这种情况下，尽管罪行是前所未有的邪恶，受害的人还活着，亲身在场；敌对派在诡计被知晓、阴谋暴露时，可以被定罪并受到惩罚。通过声明、解释和类似方式，得到了补救：因为活着的人可以为自己采取行动，死人则无法驳斥他们所承受的指控。\par

再举另一个例子。殉道者\Person{西彼廉}的书信全集通常收录在一个手抄本中。某些持有关于圣神亵渎性教义的异端者，将\Person{戴尔都良}关于天主圣三的一篇论文插入了这个集子中，这篇论文措辞有误，尽管他本人是维护我们信仰的；他们从这样制成的抄本中又抄写了许多其他抄本；并以极低的价格分发到整个庞大的\Place{君士坦丁堡}城；人们被低廉的价格所吸引，轻易地购买了这些充满隐藏陷阱却毫不知情的文献；于是异端者通过一个伟大名字的权威，为其邪恶教义找到了获取信任的手段。然而，在出版后不久，那里出现了我们的一些公教兄弟，他们能够揭露这邪恶的伪造，并尽力将尽可能多的人从错误的纠缠中召回。他们取得了部分成功。但当地仍有许多人确信圣殉道者\Person{西彼廉}持有\Person{戴尔都良}错误表达的那种信仰。\par

我再举一个文献伪造的例子。这是一个虽属近代记忆、却是原始狡诈的实例，而且胜过所有古人的故事。\par

达玛苏主教在处理阿波利纳里派回归教会的商议时，希望拥有一份阐述教会信仰的文件，由那些愿意回归者署名。他将这份文件的编纂委托给他的一位朋友，一位司铎，且是极有学问的人，通常代他处理这类事务。当他开始编写文件时，在谈论我们主的降生成人时，他感到有必要用\OriginalTerm{主的人}{Homo Dominicus}这个称呼来指称主。阿波利纳里派对这称呼感到冒犯，并开始攻击它，认为它是一个新词。于是文件的作者为自己辩护，并引用古代公教作家的权威来反驳反对者；他碰巧向一位抱怨这称呼新奇的反对者展示了一本亚大纳削主教的著作，其中出现了所讨论的词。这位收到证据的人似乎被说服了，并请求借阅这手稿，以便说服其他出于无知仍持反对的人。当他得到手稿后，他设计了一种全新的篡改方法。他先擦去包含该称呼的段落，然后又重新写上刚刚擦去的相同文字。他归还了纸张，被毫无疑问地接受了。关于这一称呼的争论再次兴起；手稿被拿了出来；所讨论的称呼被找到了，但位于有擦除痕迹的地方；而提交该手稿的人失去了所有权威，因为擦除似乎是不正当行为和篡改的证据。然而，在这个案子里，如同我之前提到的那个一样，是一个活着的人被另一个活着的人如此对待；他立即尽一切努力揭露所行的不义欺诈，将这邪恶行为的污点从无辜且从未做过此类恶事的人身上移除，并归于行为的真正作者，使他完全身败名裂。\par

既然奥力振在他的书信中亲自抱怨他曾遭受那些恶意对待他的异端者如此对待，而且类似的事在许多已故和在世的正统人士身上都发生过；既然在所引用的案例中，人们的著作被证明遭受过类似的篡改：那么，这种顽固的决心是什么，竟然拒绝在同样的情况下接纳同样的辩解，并且在情况相似时，给予一方应有的尊重，而给予另一方罪犯的耻辱？必须说出真相，不能在这一点上隐藏；因为任何人都不可能这么不公正地判断，以致对相似的情况形成不同的意见。事实上，奥力振的控告者的唆使者是那些在教会中发表冗长争论性言论，甚至写书，其全部内容皆从奥力振那里借来的人物；他们想阻止心地单纯的人阅读奥力振，唯恐他们的剽窃行为广为人知，尽管，确实，如果他们不是对老师忘恩负义，他们的借用本不会成为他们的耻辱。\par

例如，这些人中的一个，他认为自己有一种必要，如同宣讲福音一般，要在万国万民中恶言攻击奥力振；他在一大群基督徒听众面前宣称，他读过奥力振六千部著作。当然，如果他读这些著作的目的，如他惯常所宣称的，只是为了了解奥力振的过错，那么十部、二十部，至多三十部就足够了。但是读六千部书不再是希望了解这个人，而是几乎将整个人生都奉献给了他的教导和研究。那么，当他谴责那些只读过其中很少几部书的人，而这些人保持着信仰准则的神圣和虔敬的无缺时，他的话又如何值得相信呢？\par

以上所说的足以说明我们应当对奥力振的著作持何种看法。我认为，每个关心真理而非争论的人都容易同意我提出的有充分证据的论述。但如果有人执意好辩，我们却没有这样的惯例。我们中间有一个确定的惯例：在读他时，我们要持守那美好的，遵照宗徒的训令。如果我们在这些书中发现任何与公教信仰不符之处，我们怀疑那是异端者插入的，并认为它既与他的见解相异，也与我们的信仰相异。然而，如果这是我们错误的判断，我认为，我们不会因这错误而陷入危险；因为借着天主的助佑，我们自身通过避免我们所怀疑和谴责的而安然无恙；此外，我们不会在天主面前成为弟兄们的控告者（你们要记得，控告弟兄是魔鬼的特殊工作，他被称为魔鬼正是由于他是毁谤者）。而且，我们因此逃脱了针对恶言者的判决，这判决将这样的人从天主的国中分离出去。\par

% structure: p0043 source=h2 target=section reason=relative-heading-rank; base=h2

\section{奥力振\WorkTitle{\OriginalTerm{论原理}{\GreekText{Περὶ} \GreekText{᾽Αρχῶν}}}诸卷译本序}

致\Person{玛加略}，在\Place{皮内屯}，公元三九七年。\par

他很快在庞斐禄的\WorkTitle{辩护辞}之后或同时发行了\WorkTitle{\OriginalTerm{论原理}{\GreekText{Περὶ} \GreekText{᾽Αρχῶν}}}前两卷的译本。这篇序言的目的是要消除偏见，显示\Person{热罗尼莫}（虽未指名，但清晰描述）曾是\Person{鲁菲努}翻译奥力振的前驱。对\Person{热罗尼莫}的赞扬无疑是真诚的：但对他先前行动的利用却几乎无法辩护。\Person{鲁菲努}深知，\Person{热罗尼莫}对奥力振的看法已在某种程度上改变，一场关于奥力振令人不快的争论在\Place{耶路撒冷}兴起，在这场争论中，他和\Person{热罗尼莫}站在了对立的两边；因此引起的敌意虽经极大努力才得以平息，并在\Person{鲁菲努}离开\Place{巴勒斯坦}时达成了和解。这篇序言连同\WorkTitle{\OriginalTerm{论原理}{\GreekText{Περὶ} \GreekText{᾽Αρχῶν}}}的译本是\Person{鲁菲努}和\Person{热罗尼莫}之间激烈争论和最终分裂的最直接原因。\par

我知道有许多弟兄因渴求圣经知识而激动，向精通希腊文学的人请求，请他们将奥力振的作品译成拉丁文，使他成为罗马人。其中有一位是我的弟兄和同道，达玛苏斯主教曾向他提出这个请求。当他从希腊文翻译关于\WorkTitle{雅歌}的两篇讲道时，在作品前加了一篇如此优美而宏伟的序言，以致唤起了每个人阅读奥力振并热切探究其作品的渴望。他说，那位伟人的灵魂可以恰当地应用这句话：\BibleQuote{君王引我进入了内室}；并宣称奥力振在其著作中已超越所有人，而在这一作品上更超越了自己。他在此序言中所许诺的，确实是将不仅这些书卷，而且还有许多奥力振的其他作品献给拉丁读者。但我发现他如此陶醉于自己的风格，以至于追求一个更宏大的目标：他应当是书的创造者，而不仅仅是翻译者。因此，我是在继续他开创的、并由他的榜样所推荐的工作；但我无力像他那样有力而雄辩地阐述这位伟人的话语。因此我担心，由于我的过错，那个他公正地称赞为教会导师、在知识和智慧上仅次于宗徒的人，可能因我语言的贫乏而被认为地位低得多。当我想到这一点时，我倾向于保持沉默，不同意那些不断恳求我翻译的弟兄们。但是你，我最忠信的弟兄玛加略，你的影响力如此之大，以至于即使我意识到自己的不称职，也不足以使我抵抗。因此，我屈服于你的强求，尽管违背了我的决心，以便我不再受一个严厉的工头的要求；但我这样做是有条件和理解的：在翻译中，我将尽可能遵循我前辈的方法，尤其是那位我已经提到过的人的方法。他在将奥力振的七十多部他称为讲道集的书以及一些卷册翻译成拉丁文之后，继续在翻译中清除所有在希腊著作中可能引起绊脚石的原因；他这样做使得拉丁读者在其中找不到任何与我们的信仰相冲突的东西。因此，我跟随他的脚步，虽然不拥有他那样的雄辩能力，但尽可能遵循他的规则和方法，即注意不发表那些在奥力振的著作中发现的彼此不一致和矛盾的东西。关于这些差异的原因，我在潘斐路斯为奥力振著作所写的\WorkTitle{护教篇}中已经非常详细地向你说明，我还附加了一篇非常简短的论文，用我认为相当清楚的证据表明，他的著作在很多时候被异端和恶意的人篡改。尤其是你现在要求我翻译的这些书，即\OriginalTerm{论首要原理}{\GreekText{Περὶ} \GreekText{᾽Αρχῶν}}，可以译为《论首要原理》或《论掌权者》。这些书，撇开这些问题不谈，确实是极其模糊和困难的；因为他在其中讨论了哲学家们穷尽一生而毫无结果的问题。但我们这位基督徒思想家尽力将他所创造的关于创造主和受造界秩序的信条（他们曾使之服务于虚假宗教）转向健全宗教的目的。因此，每当我发现他的书中有任何与关于天主圣三的真理相悖之处——他在其他地方曾以严格正统的方式论述过——我便要么将其作为外来的、不真实的表述省略，要么用与他经常同意的信仰规则相符的措辞记录下来。无疑，有些内容他用相当晦涩的语言展开，希望迅速带过，仿佛是对那些对此类事有经验和知识的人说话；在这些情况下，我通过添加我在他其他书中读到过的更全面论述的内容来使段落清晰。我这样做是为了清晰明了；但我没有添加任何自己的东西；我只是把他自己的话还给他，尽管是从其他段落中取来的。我在序言中解释了这一点，以免那些诽谤我们的人以为他们找到了新的攻击材料。但是，让他们注意自己在乖张和好争辩中的所作所为。至于我，我承担这项费力的工作（我相信天主会应允你的祈祷而帮助我）并不是为了堵住诽谤者的口，而是为了给那些渴望真正知识的人提供知识增长的材料。我只要求每个抄写或阅读这些书的人，在天主圣父、圣子和圣神面前，并凭着我们对于来临之国的信仰、对于死者复活的保证、以及\BibleQuote{为魔鬼及其使者预备的永火}（正如他相信他不会承受那有哀号和切齿的地方作为永恒的产业，在那里\BibleQuote{火不灭，虫不死}），不要增添或删减，不要插入或更改任何写下的内容，而要将其副本与所从抄写的原本进行对照，逐字逐句地严谨校订，并且不要将未经校订或审阅的副本保留在身边，免得因为副本不清晰，造成意义的困难，在读者心中产生更大的晦涩。\par

% structure: p0047 source=h2 target=section reason=relative-heading-rank; base=h2

\section{\OriginalTerm{论首要原理}{\GreekText{Περὶ} \GreekText{᾽Αρχῶν}}第三卷序言}

\Person{鲁菲诺}此时已到了\Place{罗马}。第三卷和第四卷的翻译可能早在398年初于\Place{皮内图姆}完成。他已意识到其翻译的第一卷和第二卷所引起的强烈反响，并抱怨他作品的部分内容在尚未校正之前就被\Person{热罗尼莫}的朋友们获取，并用以诋毁他（《辩护》一18-21，二44）；但他仍继续这项工作，并为此冠以以下序言作为其正当理由。\par

读者，请在你神圣的祈祷时刻纪念我，使我堪当成为圣神的追随者。\Person{玛加略}，正是由于你的激励，我甚至可以说是受你的强迫，翻译了《\OriginalTerm{论第一原理}{\GreekText{Περὶ} \GreekText{᾽Αρχῶν}}》的前两卷。我是在四旬期完成的；那时你近在咫尺，我基督内的兄弟，你的充裕闲暇也迫使我更加勤奋。但现在你住在\Place{罗马}城另一端，我的工头较少来访，我花了更长的时间来阐述后两卷的含义。你会记得，我在前一篇序言中曾警告你，有些人发现我对\Person{奥利振}毫无诋毁时会充满愤怒：这，我想你已经发现，不久便应验了。但如果那些煽动人口出恶言的魔鬼已因那作品的第一部分而被点燃，尽管作者在其中尚未完全揭露他们的诡计，这第二部分又会造成何等影响？在这部分中，他将揭示那些潜入人心、欺骗软弱与脆弱心灵的隐秘迷宫。你会看到骚动从四面八方涌起，党派精神被激起，全城将掀起一片喧嚣，\Person{奥利振}将被传唤到法庭，因试图以福音之灯驱散无知黑暗而受谴责。但这一切对于那些努力在研习神圣事物中持守公教信仰纯正形式的人来说，实在算不了什么。\par

然而，我认为有必要提醒你我在处理先前各卷时所遵循的原则，这一原则在此也适用，即不在我的译文中放入显然与我们的信仰及作者在其他地方表达的观点相矛盾的东西，而是将其视为他人插入的、非本真的内容而略过。另一方面，无论是在前几卷还是这几卷中，我都没有省略他关于受造理性物形成的新奇观点，因为我认为信德并不主要在于这些事情，而在于他所追求的是知识与心智的运用，并且可能有一些异端需要用这种方式来回答。只是，在他在后几卷中重复其先前所说之处，为了简洁起见，我删除了某些部分。\par

那些阅读这些书的目的在于获取知识而非贬损作者的人，最好寻求比他们更有学问之人的帮助来解释它们。因为让文法学家给我们解释诗人作品的虚构和喜剧演员的可笑故事，却认为那些谈论天主、天上权能和整个宇宙，并讨论外邦哲学和异端乖谬之所有错误的书，任何人都无需教师解释就能理解，这是荒谬的。由此，人们宁愿保持无知，对困难晦涩之事轻率下判断，而不愿通过勤奋研习来理解它们。\par

% structure: p0052 source=h2 target=section reason=relative-heading-rank; base=h2

\section{\Person{阿纳斯塔西}，\Place{罗马}教会主教，致\Place{耶路撒冷}主教\Person{若望}关于\Person{鲁菲诺}品性的信}

\Person{阿纳斯塔西}致\Person{耶路撒冷的若望}的信写于401年；\Person{热罗尼莫}在写于402年上半年的《辩护》卷三第21章中提到这封信，称其为去年的信。\Person{热罗尼莫}在同一处表明，这只是\Person{阿纳斯塔西}写给东方多封相同性质信函中的一封。\Person{鲁菲诺}未曾见过此信，并拒绝相信其真实性。但似乎没有理由怀疑它。\Person{阿纳斯塔西}在\Person{亚历山大的德奥斐洛}的恳切请求下，正式谴责了奥利振主义。而\Person{鲁菲诺}翻译的\Person{奥利振}《\OriginalTerm{论第一原理}{\GreekText{Περὶ} \GreekText{᾽Αρχῶν}}》和\Person{庞斐禄}的《为奥利振辩护》，连同他自己论奥利振著作被篡改的书，在\Place{罗马}被视为对奥利振主义的一般辩护。然而，\Person{鲁菲诺}不断上诉，特别是在他给\Person{阿纳斯塔西}的《辩护》中，诉诸他曾在其中领受圣职的\Place{耶路撒冷}教会。他说，我的信德就是\Place{耶路撒冷}所宣讲的信德。因此，\Person{阿纳斯塔西}谴责奥利振就会被理解为谴责\Person{鲁菲诺}，并可能似乎谴责他的主教\Person{耶路撒冷的若望}。这可以说明信函开头过分恭维的原因。此外，\Person{若望}曾写信给\Person{阿纳斯塔西}咨询关于\Person{鲁菲诺}的事，这可能意味着某种有利于\Person{鲁菲诺}的行动；但\Person{热罗尼莫}知道信的内容而\Person{鲁菲诺}不知道，这一事实似乎表明\Person{若望}主教与\Person{热罗尼莫}的关系变得更友好，而与\Person{鲁菲诺}则疏远了。\par

1. 你对我亲爱的主教所表达的赞许之言，正是你长期考验之爱的新标志。这是你赐予我的高度赞扬，是对我辛劳最慷慨的认同。我为这份爱的证明感谢你；在我卑微中远远跟随你，我用言语的贡品来尊崇你的圣德的光辉，以及主赏赐给你的诸般德行。你远超众人，你的赞美之光如此显赫，以至于我所能用的言语无法相称你的功德。然而，你的荣耀在我心中激起如此敬佩，以致我无法放弃描述它的尝试，尽管我永远无法充分做到。首先，你从你伟大精神的宁静天空赐予我的赞美，构成了你自己荣耀的一部分：因为是你主教职的威严，如同太阳照耀世界相对的一方，将其光辉反射到我们身上。你毫无保留地赐予我友谊；你不以批评的天平衡量我。如果赞美我是对的，难道你的赞美不应该回响给你吗？因此，我恳求你，为了你也不亚于为了我，不要再当面赞美我。我要求这一点有两个原因：如果赞美是不当的，它必定在你的主教兄弟心中激起痛苦；如果是真实的，它必定使他脸红。\par

2. 让我谈谈你信的主题。关于\Person{Rufinus}，你曾屈尊向我请教意见，他必须将自己的良心带到天主威严的审判台前。他必须自己想方设法，看如何能在天主面前证明自己保持了对祂的真诚忠诚。\par

3. 至于\Person{Origen}，他将此人著作译成我们语言，我从前既不知道他是谁，或者他如何表达自己的思想，现在也不想了解。但是关于这件事在我心中留下的感受，我愿与你的圣德稍作交谈。我所获得的印象是这样的——这已通过我们城里的人诵读\Person{Origen}著作的部分，以及它投在他们身上的那种盲目迷雾，清楚地显明出来——他的目的是要瓦解我们的信德，即宗徒的信德，并已由教父的传统所确认，通过将我们引入弯曲的道路。\par

4. 我想知道，将这部著作译成拉丁语是什么用意。如果译者借此意在指出作者的错误，并向世界宣告他的可憎行径，那很好。那样，他将使一个长期背负舆论不利重压的人遭受应得的仇恨。但如果通过翻译所有这些邪恶之事，他意在表示赞同，并以此意让世界阅读，那么他付出如此巨大辛劳所建造的大厦，除了使罪过成为他自己意志的行为，并为他意外的支持赋予认可，从而颠覆自宗徒时代至今天主教基督徒所持守的真信德中所有首要之事外，别无他用。\par

5. 愿这样的教导远离罗马教会的天主教体系。我们绝不能以任何可能的方式，接受我们作为法律和正义之事所谴责的东西为合理。因此，我们确信，基督我们的天主，其眷顾遍及整个世界，当祂说我们完全不可能接纳那些玷污教会、颠覆其久经验证的道德体系、冒犯所有目睹我们行为之人的耳朵、为纷争、愤怒和分裂奠定基础的教义时，祂对我们表示赞同。这就是促使我写信给我们在主教职中的兄弟\Person{Venerius}的动机，这封信的性质，尽管出于我的软弱，却以极大的谨慎和勤奋写成，你将从我现在所附的内容中认识到这一点：那么，由此，翻译这部著作的人从何处获得并保持这种无罪的确信，我并不急于知道：它并未引起我无谓的惊慌。我当然不会遗漏任何能使我守护我所辖人民中的福音信德的事，并尽我所能，警告那些组成我身体的一部分、无论他们生活在世界何处之人，不要让任何对亵圣作者的翻译渗透并从中兴起，这些翻译将试图通过传播其自身的黑暗来动摇虔敬之人的心灵。此外，我不能缄默一件事，它给了我极大的喜悦，即我们皇帝颁布的法令，其中警告每一位事奉天主的人不要阅读\Person{Origen}，所有被证实阅读他亵圣著作的人都被皇帝判决定罪。我用这些话表达了我的正式裁决。\par

6. 你因人们对我们对待\Person{Rufinus}方式的抱怨而感到困扰，以致你对某些人怀有模糊的怀疑。但我将用取自圣经的一个例子来回应你的这种感觉，即经上所说：\BibleQuote{人看人不像天主看人；因为天主看人心，人只看外貌。}因此，我亲爱的弟兄，放下你所有的偏见。用你自己不偏不倚的判断衡量\Person{Rufinus}的行为；问你自己，他是否没有将\Person{Origen}的话译成拉丁文并加以认可，以及一个鼓励他人所犯恶行的人，与犯罪者本人有何不同。无论如何，我恳请你确信这一点：他完全与我们无关，我既不知道也不想知道他在做什么或住在哪里。我只补充一点：他应当考虑在哪里可以获得赦免。\par

% structure: p0060 source=h2 target=section reason=relative-heading-rank; base=h2

\section{\WorkTitle{息斯笃语录}翻译序言}

约公元三〇七年写于\Place{Aquileia}。\par

\Signature{鲁菲诺致友人阿普罗尼亚诺}\par

我知道，正如羊群听见自己的牧人呼唤便欣然前来，同样，在信仰之事上，人们最乐意聆听一位用他们自己的语言讲话的教师的劝诫。因此，我亲爱的\Person{Apronianus}，当那位虔诚的女士——她是我在基督里的女儿，如今却是你在基督里的姊妹——命令我为她撰写一部无需多大努力便能理解的论著时，我便以一种非常直白、简明的风格，将\Person{Xystus}（据说此人就是在罗马被称为\Person{Sixtus}的那位，他获得了主教兼殉道者的荣耀）的著作译成了拉丁文。我想，当她阅读时，她会发现这部著作表达得如此简洁，以至于每一行都展开宏大的含义；如此有力，以至于仅仅一行长的句子就足以作为一生的训练；却又如此简单，以至于一个在她阅读时从她肩后窥看的人可能会怀疑我是否智力低下。整部作品如此简明，以至于她可能永远不愿放手。整本书几乎不比我们先人的一枚指环更大。确实，一个藉着天主圣言学会将世界的装饰视为渣滓的人，如今从我手中接受一条由圣言和智慧制成的项链作为装饰，似乎是合宜的。目前，就让这本小书像一枚指环一样常握在手中；但不久它就会进入宝库，完全藏在心里，并从其内室产生教导和参与一切善工的种子。我又加上了几段一位虔诚父亲对儿子所说的精选格言，但都非常简短，以至于这整部小作在希腊文中可以恰当地称为\OriginalTerm{手册}{Enchiridion}，或在拉丁文中称为\OriginalTerm{指环}{Annulus}。\par

% structure: p0064 source=h2 target=section reason=relative-heading-rank; base=h2

\section{为\WorkTitle{教会史两卷}所写的序言，由\Person{Rufinus}附于其\Person{Eusebius}译本之后}

致\Place{Aquileia}主教\Person{Chromatius}，公元401年。\par

\EditorialNote{关于写作的缘由和日期，见导言，第412页。}\par

据说，技艺高超的医生有这样的习惯：当他们察觉到某一种流行病临近我们某一座城市时，便提供某种药物，无论是固体的还是液体的，人们可以将其用作预防，以保护自己免受悬临的毁灭之害。您效法了医生的这种方法，我尊敬的神父\Person{Chromatius}，当意大利的大门被哥特人的统帅\Person{Alaric}攻破，一场疾病和瘟疫便如此倾泻在我们身上，蹂躏了遍地的田野、牲畜和人民。那时，您寻求对付残酷和毁灭的补救方法，使那些衰弱的心灵得以从流行病的传染中抽身，并通过关注更高尚的事业保持平衡。您通过命令我将那位最有学问的\Place{Cæsarea}的\Person{Eusebius}用希腊文写成的教会史译为拉丁文，实现了这一点。您认为，听到诵读的人的心灵可能会被它紧紧抓住，以至于在对过往事件知识的渴望中，他们可能在某种程度上忘却了当前的苦难。我试图推辞这项任务，因为我的软弱无法胜任，并且年复一年，我已经失去了拉丁语的使用能力。但我想，您的命令不应与您在宗徒秩序中的地位相违背。因为，当旷野中的群众饥饿时，主对他的宗徒们说：\BibleQuote{你们给他们吃吧。} 其中的一位宗徒\Person{Philip}，没有拿出藏在宗徒行囊里的饼，却说那里有一个小孩，带着五个饼和两条鱼。（参\BibleRef{若 6:9}；\BibleRef{玛 14:16}；\BibleRef{谷 6:37}；\BibleRef{路 9:13}）他知道，如果在实现过程中使用了某个小孩的服务，神圣德能的彰显并不会因此减损。他谦卑地为自己的行动辩解，补充说：\BibleQuote{这为这许多人还算什么？} 这样，神圣的力量通过其作用的艰难和绝望处境而更加明显。我感到，既然您是宗徒秩序的后裔，您可能记起了\Person{Philip}的榜样，当您看到喂饱群众的时刻到来时，您找来了一个小男孩的服务，他可能贡献出他所领受的五个饼——双倍计数——并且，为了成全福音的预象，他还可能加上他自己捕获的两条小鱼。因此，我尝试执行您所命令的事，确信我的缺乏经验会因发令者的权威而得到谅解。\par

我必须指出我在处理这部作品的第十卷时所采取的方法。在希腊原文中，它与事件进程关系不大。除了一小部分外，大部分内容都是讨论，旨在赞扬某些主教，而丝毫无助于我们对事实的了解。因此，我删掉了所有这些多余的材料；凡属于真实历史的内容，我都加到了第九卷中，我使他的历史以第九卷结束。第十和第十一卷是我自己编纂的，一部分来自前一代的传统，一部分来自我记忆中的事实；我将这些加到前面的卷册中，就像那两条鱼加在饼上一样。如果您赐予它们您的认可和祝福，我将确信它们足以满足群众。现在完成的作品包含了从救主升天至今的事件；我自己的两卷则包含了从\Person{Constantine}时代（当时迫害结束）直到\Person{Theodosius}皇帝去世期间的事件。\par

以下注释出现在\Person{Rufinus}拉丁文版的\Person{Eusebius}著作第九卷末尾。\par

至此，\Person{Eusebius}为我们提供了历史的记录。至于随后的事件，直到现在，正如我从上一代人的著作中发现的，或依据我自己的知识，我将把它们添加上去，在此事上也尽我所能，服从我们在天主内的父亲的命令。\par

% structure: p0071 source=h2 target=section reason=relative-heading-rank; base=h2

\section{\Person{鲁菲诺}翻译\Person{奥利振}\WorkTitle{\BibleBook{圣咏}{}第三十六、三十七、三十八篇注释}序言}

致\Person{阿普洛尼亚诺}，或在\Place{罗马}，或在\Place{阿奎莱亚}，约公元398年至407年之间。\par

对第三十六、三十七和三十八篇圣咏的全部阐述具有伦理性质，旨在强制推行更正确的生活方式；时而教导皈依与补赎之道，时而教导净化与进步之道。因此，我认为最好将其译为拉丁文献给你，我最亲爱的儿子阿普洛尼亚诺，首先将其编排为九篇短讲道（希腊语称为\OriginalTerm{讲道}{Homilies}），并整合为一个整体；这样，这篇旨在从各方面纠正和提升道德生活的论述，便汇集于单卷之中。我的翻译至少会有所帮助，使读者毫不费力地掌握作者的用意，此处已充分揭示，并将他所要求的简单生活以清晰的思想和朴素的言语传达给他；这样，预言的声音不仅可达到男人，也可达到敬畏天主的妇女，并为简单人的心智增添精致。但我担心那位虔诚的女士，她是我的女儿，但在基督内是你的姊妹，可能会认为她不必为我的工作感谢我，如果它带给她的只是令人困惑的思想和棘手的问题：因为人的身体几乎无法维持，如果天主眷顾只以骨骼和肌肉构成它，而没有将柔软组织的轻松和优雅与之融合。\par

% structure: p0074 source=h2 target=section reason=relative-heading-rank; base=h2

\section{\Person{鲁菲诺}翻译\Person{奥利振}\WorkTitle{\BibleBook{罗马}{书}注释}序言}

致\Person{赫拉克利}，在\Place{阿奎莱亚}，约公元407年。\par

我本意是要驾着我航行的小舟，贴着宁静之地的海岸，从希腊的水池中钓出几条小鱼；但你却强迫我，\Person{赫拉克利}兄弟，让我扬帆出海，驶向深海；你说服我放下手头正在翻译的\Quote{铁人}晚年所写讲道集的工作，为你打开他论述\BibleBook{罗马}{书}的十五卷著作。在这些书中，当他试图表达宗徒的思想时，他被卷入如此深的大海，以至于跟随他的人可能会担心，在思想之伟大中如同在浪涛之浩瀚中淹死。此外，你也没有考虑到这一点：我的气息不足以吹响像他那样雄辩的宏大号角。而超越所有这些困难的是，这些书籍本身已被篡改。几乎在所有图书馆中（我承认没人知道这是怎么发生的），有些卷册从著作中缺失了；而要补充它们，并在拉丁译本中恢复工作的连续性，超出了我的才能，但正如你在提出要求时必须知道的那样，这是天主的特别恩赐。然而，你还加上了一句，以便我所承担的工作无所欠缺：我最好将这十五卷的全部内容（希腊文长达四万行或更多）加以缩写，并将其压缩到适中的篇幅。你的指示确实艰难，可能会被认为是由一个不考虑此类工作负担的人强加的。不过，我将尝试，希望藉着你的祈祷和主的恩惠，在人看来不可能的事可能成为可能。但我们现在，如果你愿意，请听\Person{奥利振}本人为他所从事的工作所写的序言。\par

% structure: p0077 source=h2 target=section reason=relative-heading-rank; base=h2

\section{\Person{鲁菲诺}附于其翻译的\Person{奥利振}\WorkTitle{\BibleBook{罗马}{书}注释}之后的结束语}

致\Person{赫拉克利}，在\Place{阿奎莱亚}，可能约407年。\par

现在，我相信，对\WorkTitle{罗马书}注释的工作已经令人满意地完成了，这部注释的写作耗费了极大的劳力和时间。我承认，我最亲爱的兄弟\Person{赫拉克利乌斯}，在试图回应你的请求时，我几乎忘记了那诫命：\Quote{不要承担超过你力量的负担}。即使在奥利振其他著作的拉丁文翻译中——这些翻译是应你恳切的要求，更确切地说，是作为一项工匠任务强加于我的——劳力也是巨大的；因为我以补充\Person{奥利振}在教会讲堂中即兴演讲的内容为目标；因为他在那里的目的是为了教诲而应用主题，而非解释经文。在\WorkTitle{讲道集}和关于\WorkTitle{创世纪}与\WorkTitle{出谷纪}的简短讲座中，尤其是关于\WorkTitle{肋未纪}的讲座中，我这样做了，他在那里以劝诫的方式说话，而我的翻译则采取了论述的形式。我承担起这个补充所需的责任，是因为我认为他经常在讲道式说话中采用的那种提出问题然后留而不决的做法，可能会使拉丁读者感到不快。关于\WorkTitle{耶稣·纳维}和\WorkTitle{民长纪}以及\WorkTitle{圣咏集 第三十六、三十七、三十八篇}的作品，我简单地按原样翻译，没有花费太大劳力。然而，在其他我上面提到的情况下，我花费了大量劳力来补充\Person{奥利振}所省略的内容，而在\WorkTitle{罗马书}这部作品中，由于前言中所述的原因，落在我身上的劳力是巨大而复杂的。但若没有其他情况中恶意的心态以侮辱回报我的辛劳和守夜，并且我的研究没有招致诽谤的回报，我的辛劳没有招致毁掉我的阴谋，那么这些辛劳将只有快乐。因为对付这些人，我要经历一种新的指控。他们对我说：当你写这些东西时，其中发现许多篇章节是出于你本人的创作，你应该把你的名字放在标题中，让它这样写：\Quote{\Person{鲁菲努斯}关于（例如）\WorkTitle{罗马书}的注释之书}；因为，他们说，在外教作家的情况下，标题中的名字不是被翻译的希腊作者的名字，而是翻译他的拉丁作者的名字。但这种以归功于我作品为形式的殷勤，并非出于对我的爱，而是出于对那位作者的恨。我更加注意我的良心，而不是我的名声；可能显而易见的是，我添加了一些东西来补充所缺，并且缩短了过于冗长的部分；但窃取那个奠定建筑基础之人的标题，我认为是不对的。当然，我承认，读者在检查作品后，有权将其归于他认为合适的人；但我的意图不是寻求学生的掌声，而是那些希望得到教诲之人的益处。\par

我接下来将转向那个很久以前就交给我、但现在由主教\Person{高登提乌斯}更加迫切要求的工作，即将被称为\WorkTitle{克肋孟的明辨}的书翻译成拉丁文，\Person{克肋孟}是罗马的主教，宗徒的继承者和同伴。在这项工作中，我很清楚，按照一般规则，我将有劳作叠加劳作。在这种情况下，我会做朋友们希望的事，我会把自己的名字放在作品的标题中，虽然我也会放作者的名字。它将被叫做\WorkTitle{鲁菲努斯的克肋孟}。如果主使我能够完成这项任务，之后我将回到你所期望的，并在天主的助佑下，对\WorkTitle{户籍纪}或\WorkTitle{申命纪}（因为这是我关于\WorkTitle{梅瑟七书}整个工作中唯一缺少的）说些什么；或者，在主引导下，我尽可能在\Person{保禄}宗徒其余的书信上写下一些东西。\par

% structure: p0081 source=h2 target=section reason=relative-heading-rank; base=h2

\section{\Person{奥利振}《\WorkTitle{户籍纪讲道集}》序言}

致\Person{乌尔萨基乌斯}。写于410年。\par

我亲爱的弟兄，我本可以用蒙福的老师的话来称呼你：\Quote{亲爱的多纳图斯，你提醒我这件事很好，因为我清楚记得我的承诺——要将阿达曼提乌斯晚年所写关于梅瑟法律的一切收集起来，并为我们的同胞译成拉丁语。但如他所说，那时节并不适宜履行我的承诺，反而充满了风暴与混乱。}当一个人惧怕敌人的投射物，当他眼前是城市与乡村的荒芜，当他不得不逃避海上的危险，连流亡中也没有安全时，笔怎能自如地移动呢？正如你亲眼所见，蛮族已近在咫尺；他纵火焚烧了\Place{雷吉乌姆}城，我们唯一的屏障就是那分隔\Place{意大利}与\Place{西西里}的极狭窄的海域。在这样的处境中，哪里还有闲暇写作，尤其是翻译一部作品——其职责不是发挥自己的见解，而是表达别人的意思呢？然而，当宁静的夜晚来临，我们的心神从敌人进攻的恐惧中解脱，获得了至少些许思考的闲暇，我便着手工作，作为困难中的慰藉，减轻我们旅途的重担，将\Person{奥利振}（即\Person{阿达曼提乌斯}）关于\BibleBook{}{户籍纪}所写的一切——无论是讲道，还是称为\WorkTitle{选录}的作品——收集整理，译为拉丁文。\Person{乌尔萨齐乌斯}，是你敦促我做这事，并全力相助；你如此急切，竟觉得担任秘书的青年动作太慢。然而，弟兄，我想向你指出：这种研读圣经的方法，其目标并非像你在注释书中看到的那样逐句处理，而是为理解开辟一条道路，使读者不至于懈怠，而是如经上所记，激发自己的心神，引申出其含义，并在听到善言之后，以自己的智慧增添其内容。我以此方式努力提供了你所渴望的全部解释。如今，我所发现的关于法律的著作中，唯独缺少对\BibleBook{}{申命纪}的简短注释；若天主愿意，且恢复我的视力，我希望将这些补充到著作整体中。事实上，我至爱的儿子\Person{皮尼亚努斯}——因其贞洁的言谈而令我喜乐，我在流亡中与他真心基督徒般的陪伴同行——还向我要求其他工作。但请你和他一同祈祷：愿主与我们同在，在我们的时代赐下平安，怜悯在困苦中的人，并使我们的工作结出果实，以建树读者。\par

% structure: p0084 source=h2 target=section reason=relative-heading-rank; base=h2

\section{\Person{鲁菲诺}序言}

我知道，许多弟兄因渴望认识圣经，曾请求一些精通希腊学问的杰出人士将\Person{奥利振}译成拉丁文，以便\Place{罗马}读者能阅读。其中，我们的兄弟和同僚在\Person{达玛苏斯}主教的恳切请求下，将\WorkTitle{雅歌}的两篇讲道从希腊文译成拉丁文，并为该作品撰写了如此优雅而高贵的序言，以致激发了每个人对阅读和研究\Person{奥利振}的热切渴望，他引用说：\Quote{\BibleQuote{君王引我进入了他的内室}}\BibleRef{歌 1:4}，并宣称\Person{奥利振}在其他作品中超越所有作者，而在\WorkTitle{雅歌}中他甚至超越了自我。的确，他在那篇序言中承诺，将把关于\WorkTitle{雅歌}的书籍以及\Person{奥利振}的其他许多作品以拉丁译本呈现给\Place{罗马}读者。但他更喜欢自己的创作，追求更享盛誉的目标，即成为作品的作者而非译者。因此，我们着手进行他这样开始并认可的这项事业，尽管我们无法以与他那样卓越的口才相媲美的优雅风格来写作；所以我担心，由于我的过失，会导致这样的结果：那位他理应尊为继宗徒之后教会中第二位知识与智慧之师的人，会因我语言的贫乏而显得远非其本相。这一考虑时常回响在我心中，使我保持沉默，未能应允弟兄们的众多恳求，直到您，我最忠实的兄弟\Person{玛加略}，您如此巨大的影响力使我无法再以我的不熟练来抵抗。因此，为了不让您成为过于严厉的索求者，我违背自己的决心而让步了；但条件是，在我的翻译中，我将尽可能遵循我的前辈们所遵守的规则，尤其是上面提到的那位杰出人物，他在将\Person{奥利振}那些被称为\Emph{讲道集}的七十多篇论文以及相当多的关于宗徒的著作译成拉丁文后（这些著作的希腊原文中有许多绊脚石），在他的译文中如此润色和修正，以致拉丁读者不会遇到任何与我们的信仰不一致之处。因此，我们尽我们所能效法他的榜样；即使没有相等的口才，至少以同样的严格规则，注意不重现\Person{奥利振}著作中那些不一致且相互矛盾的表述。这些变异的原因，我们在\Person{庞斐禄}为辩护\Person{奥利振}著作所写的\WorkTitle{护教篇}中更为自由地解释了，我们在其中附加了一篇短论，以我认为确凿的证据表明，他的书籍在许多地方被异端者和心怀恶意的人篡改，尤其是那些您现在要求我着手翻译的书籍，即那些可题为\WorkTitle{论首要原理}或\WorkTitle{论元首}的书卷，这些书卷在其他方面也确实充满晦涩和困难。因为他讨论了那些哲学家穷尽一生也未能发现的课题。但在这里，我们的作者尽其所能，努力将对造物主的信仰和受造物的理性本性（后者曾被用于邪恶目的）转向为宗教服务。因此，如果我们在他的著作中发现任何与他本人在其他作品中关于\OriginalTerm{天主圣三}{Trinitas}的虔诚观点相悖的陈述，我们或是将其视为已被篡改而非\Person{奥利振}所作而省略，或是根据我们在他本人经常确认的规则来呈现它。确实，如果他为了快速推进，有时对精通知识的人说话时表达得隐晦，为了使文句更清晰，我们添加了他在其他作品中关于同一主题的更完整表述，以解释为目的，但未添加我们自己的东西，只是将他自己的（尽管出现在他作品的其他部分）归还给他。\par

因此，我在序言中说了这些劝诫的话，以免诽谤者或许认为他们再次发现了指控的借口。但让乖僻好辩的人小心他们在做什么。因为我们暂时承担了这一繁重的工作，如果天主题您的祈祷应允，并不是为了堵住诽谤者的口（这是不可能的，尽管天主或许会这样做），而是为了给那些渴望在这些知识中进步的人提供材料。确实，在圣父、圣子、圣神面前，我恳求和劝诫每一个可能抄写或阅读这些书籍的人，凭他们对来世天国的信仰、死者复活的奥迹，以及为魔鬼及其使者预备的永火\BibleRef{玛 25:41}，好使他不永远承受那\BibleQuote{有哀哭和切齿的地方}\BibleRef{玛 8:12}，以及那里的\BibleQuote{火不灭、虫不死}\BibleRef{谷 9:48}，他\BibleQuote{不为圣经增加什么，也不从中删减什么}\BibleRef{申 4:2}，不作添加或改动，而是将自己的抄本与据以抄写的副本核对，按字形订正和区分，使他的手稿没有错误或模糊，免得因抄本模糊而造成的理解困难给读者带来更大困扰。\par

\SourceRecord{Translated by W.H. Fremantle.；Nicene and Post-Nicene Fathers, Second Series；Vol. 3.；Edited by Philip Schaff and Henry Wace.；Buffalo, NY: Christian Literature Publishing Co.,；1892.；<http://www.newadvent.org/fathers/2712.htm>.}
